\chapter{Der Schaltplan}
\label{kap:editor-schaltplan}

\section{Grundsätzliches: Plan und Text}

Der Text der Datei ist maßgeblich; der Plan ist eine zweite Ansicht davon.
Nach jeder Textänderung liest der Editor die Datei neu und baut den Plan
auf: Bauteile und ihre Knoten kommen aus \code{"components"}, Lage,
Leitungen, Labels und Masse aus \code{"schematic"}. Umgekehrt schreibt jede
Änderung im Plan nur die betroffenen Stellen des Texts – Name und
\code{"nodes"} geänderter Bauteile, neue und gelöschte Bauteile und den
Abschnitt \code{"schematic"}. Kommentare, Reihenfolge und Formatierung der
Datei bleiben erhalten, und jede Änderung im Plan ist genau ein Schritt der
gemeinsamen Undo-Geschichte.

Was sich nicht zeichnen lässt – Gleichungen, Ereignisse, Analysen,
Optionen – bleibt im Text. Bauteile ohne Symbol im Plan (etwa die Kopplung
\code{K}) erscheinen trotzdem und lassen sich wie alle anderen auswählen.

\begin{hinweis}[Dateien ohne Plan]
  Öffnet man eine Datei ohne \code{"schematic"}, ordnet der Editor die
  Bauteile in Zeilen an und setzt an jeden Anschluss ein Netzlabel mit dem
  Knotennamen aus dem Text. Die Schaltung ist damit vollständig und richtig
  verbunden, nur noch nicht schön gezeichnet. Sobald man im Plan etwas
  ändert, schreibt der Editor den Abschnitt \code{"schematic"} in die Datei.
\end{hinweis}

\section{Ansicht: verschieben und zoomen}

\begin{center}
\begin{tabularx}{\textwidth}{@{}P{5.2cm}L@{}}
  \toprule
  Bedienung & Wirkung \\
  \midrule
  Mausrad & Zoomt um den Mauszeiger. \\
  \taste{Strg} + Mausrad & Zoomt (auch auf dem Touchpad). \\
  Touchpad: wischen mit zwei Fingern & Verschiebt die Ansicht. \\
  Touchpad: spreizen & Zoomt. \\
  Touchpad: zweimal tippen mit zwei Fingern (macOS) & Zeigt alles. \\
  mittlere Maustaste ziehen & Verschiebt die Ansicht. \\
  \taste{Leertaste} halten und ziehen & Verschiebt die Ansicht. \\
  \taste{F} & Zeigt die ganze Schaltung. \\
  \bottomrule
\end{tabularx}
\end{center}

Fährt die Maus über einen Anschluss oder eine Leitung, nennt ein Tooltip
den Knoten (»Knoten „out“«, »Masse«) und nach einer Simulation seine
Spannung im Arbeitspunkt. Alle Leitungen desselben Netzes werden dabei
hervorgehoben.

\section{Auswählen}

\begin{itemize}
  \item \textbf{Klick} auf ein Bauteil, eine Leitung, ein Label oder eine
    Masse wählt es aus; ein Klick ins Leere hebt die Auswahl auf.
  \item \textbf{Umschalt-Klick} fügt hinzu oder nimmt wieder heraus.
  \item \textbf{Rechteck aufziehen} (im Leeren beginnen) wählt alles darin;
    mit \taste{Umschalt} kommt es zur bisherigen Auswahl hinzu.
\end{itemize}

Ein einzelnes ausgewähltes Bauteil wird im Text markiert, und
\menue{Eigenschaften} zeigt sein Formular. Umgekehrt wählt der Plan das
Bauteil aus, in dessen Text der Cursor steht.

\section{Bauteile platzieren}

Bauteile kommen aus der Liste \menue{Bauteile}:
\begin{itemize}
  \item \taste{A} im Plan setzt den Cursor ins Suchfeld der Liste. Ein paar
    Buchstaben genügen (»kon«, »npn«, »puls«); die Eingabetaste nimmt den
    ersten Treffer.
  \item Ein Klick auf einen Eintrag der Liste wählt ihn ebenso.
\end{itemize}

Danach hängt das Bauteil am Mauszeiger. \taste{R} dreht es um 90° im
Uhrzeigersinn, \taste{X} spiegelt es; ein Klick setzt es ab. Das Werkzeug
bleibt aktiv, damit man gleich weitere setzen kann – jedes bekommt den
nächsten freien Namen (\code{R1}, \code{R2}, \dots). \taste{Esc} oder ein
Rechtsklick beendet das Platzieren.

Die Liste enthält:

\begin{center}
\begin{tabularx}{\textwidth}{@{}lL@{}}
  \toprule
  Eintrag & Vorgabe im Text \\
  \midrule
  Widerstand, Kondensator, Spule & \code{"value": "1k"}, \code{"1u"}, \code{"1m"} \\
  Spannungsquelle & \code{"dc": 1} \\
  Sinusquelle & \code{"dc": 0, "ac": 1}, Sinus 1\,V, 1\,kHz \\
  Pulsquelle & Puls 0 \dots 1\,V, Periode 1\,ms \\
  Stromquelle & \code{"dc": "1m"} \\
  Diode, Ideale Diode & ohne Modell bzw. \code{"vf": 0.7} \\
  NPN, PNP, N-MOSFET, P-MOSFET, N-JFET, P-JFET & ohne Modell (Standardwerte von SPICE) \\
  Spannungs-/stromgesteuerte Quellen E, G, F, H & Verstärkung 1 bzw. \code{"control": "V1"} \\
  Verhaltensquelle, Spannung / Strom & \code{"v": "0"} bzw. \code{"i": "0"} \\
  Schalter, spannungsgesteuert & \code{"vt": 0.5} \\
  Leitung & \code{"z0": 50, "td": "1n"} \\
  Kopplung zweier Spulen & \code{"inductors": ["L1", "L2"], "k": 0.99} \\
  Teilschaltung »\textit{name}« & je Teilschaltung der Datei ein Eintrag \\
  Modelle der Bibliothek & beim Suchen, z.\,B. »NPN\_GEN« (Kapitel~\ref{kap:editor-bibliothek}) \\
  Modellbibliothek \dots, Neue Teilschaltung \dots & öffnen die jeweiligen Dialoge \\
  \bottomrule
\end{tabularx}
\end{center}

Die Werte sind nur Vorgaben; sie werden in den Eigenschaften oder im Text
geändert. Ein Transistor ohne \code{"model"} rechnet mit den
Standardparametern von SPICE – sinnvoller ist meist ein Modell aus der
Bibliothek.

\section{Leitungen}

\taste{W} startet eine Leitung, ebenso der Eintrag \menue{Leitung zeichnen}.

\begin{itemize}
  \item Ein Klick setzt den Anfangspunkt, jeder weitere eine Ecke. Zwischen
    zwei Klicks verläuft die Leitung rechtwinklig; \taste{/} wechselt, ob
    sie zuerst waagrecht oder zuerst senkrecht geht.
  \item Trifft ein Klick einen Anschluss oder eine bestehende Leitung, endet
    die Leitung dort von selbst.
  \item Doppelklick, \taste{Esc} oder die Eingabetaste beenden die Leitung
    an der letzten Ecke; \taste{Entf} nimmt die letzte Ecke zurück; ein
    Rechtsklick bricht ab.
\end{itemize}

\paragraph{Was verbunden ist.}
Ein Leitungsende oder eine Ecke, die auf einer anderen Leitung liegt,
verbindet beide (T-Abzweig, mit einem Punkt gezeichnet). Zwei Leitungen, die
sich nur kreuzen, sind \emph{nicht} verbunden. Ein Anschluss ist verbunden,
wenn eine Leitung an ihm endet oder über ihn hinweg läuft.

\section{Masse, Netzlabels und Ports}

\begin{description}[style=nextline,leftmargin=2em]
  \item[Masse (\taste{G})]
    Jedes Netz mit einem Massesymbol ist der Knoten \code{"0"}.
  \item[Netzlabel (\taste{N})]
    Ein Label gibt seinem Netz einen Namen. Alle Labels mit demselben Namen
    sind verbunden, auch ohne Leitung dazwischen – so lassen sich
    Versorgungen und Signale ohne lange Leitungen verteilen. Doppelklick
    auf ein Label (oder \menue{Umbenennen \dots} im Kontextmenü) ändert den
    Namen; die Knoten aller angeschlossenen Bauteile im Text ändern sich mit.
  \item[Port (\taste{P})]
    Nur im Plan einer Teilschaltung: die Verbindung nach außen
    (Kapitel~\ref{kap:editor-teilschaltungen}).
\end{description}

\section{Bearbeiten}

\begin{center}
\begin{tabularx}{\textwidth}{@{}P{4.2cm}L@{}}
  \toprule
  Taste & Wirkung auf die Auswahl \\
  \midrule
  Ziehen mit der Maus & Verschieben; Leitungen an den Anschlüssen gehen mit und bleiben rechtwinklig. \\
  Pfeiltasten & Um ein Rasterfeld verschieben, mit \taste{Umschalt} um fünf. \\
  \taste{R} & Drehen um 90° im Uhrzeigersinn. \\
  \taste{X} & Spiegeln. \\
  \tasten{Strg}{D} & Verdoppeln (mit neuen Namen). \\
  \taste{Entf} & Löschen (Bauteile auch aus dem Text). \\
  \tasten{Strg}{G} & In eine Teilschaltung umwandeln (Abschnitt~\ref{sec:auslagern}). \\
  Doppelklick auf Bauteil & Eigenschaften zeigen; bei einer Teilschaltung diese öffnen. \\
  \tasten{Alt}{Klick} & Spannung der Leitung bzw. Strom des Bauteils als Kurve. \\
  \bottomrule
\end{tabularx}
\end{center}

Ein \textbf{Rechtsklick} öffnet ein Kontextmenü mit allen Befehlen, die an
dieser Stelle möglich sind – mit ihren Tasten: Eigenschaften, Strom oder
Spannung zeigen, Teilschaltung öffnen, Umbenennen, Drehen, Spiegeln,
Verdoppeln, Löschen, In Teilschaltung umwandeln, Leitung zeichnen, Masse,
Netzlabel, Port, Bauteil, Alles zeigen und gegebenenfalls Zur
Hauptschaltung.

Die Hinweiszeile unter dem Plan nennt immer die Tasten des aktiven
Werkzeugs.

\section{Wie Knotennamen entstehen}

Aus dem Plan bestimmt der Editor, welche Anschlüsse ein gemeinsames Netz
bilden, und gibt jedem Netz einen Namen:

\begin{enumerate}
  \item Ein Netz mit Massesymbol heißt \code{"0"}.
  \item Ein Netz mit einem Label trägt dessen Namen.
  \item Sonst behält es den Namen, den seine Anschlüsse bisher im Text
    hatten. Ausgaben wie \code{v(out)} bleiben deshalb gültig, auch wenn man
    Bauteile verschiebt oder neu verdrahtet.
  \item Neue Netze ohne Namen heißen \code{n1}, \code{n2}, \dots
\end{enumerate}

Ein Anschluss, an dem nichts hängt, bekommt einen eigenen Knoten. Der
Simulator wird dann meist melden, dass dieser Knoten keinen Gleichstrompfad
nach Masse hat (Kapitel~\ref{kap:fehlersuche}).

\section{Wenn Plan und Text sich widersprechen}

Weil beide bearbeitbar sind, können sie auseinanderlaufen: Ändert man im
Text die Knoten eines Bauteils, ohne die Leitungen im Plan anzupassen,
verbindet der Text etwas anderes als die Zeichnung. Der Editor zeigt das als
Warnung in \menue{Probleme} (»Schaltplan: R1, Anschluss 2 liegt im Plan an
„n1“, im Text an „out“«) und mit einem roten Ring am Anschluss. Gerechnet
wird immer, was im Text steht. Die nächste Änderung im Plan schreibt die
Knoten des Plans in den Text und beseitigt den Widerspruch.

Enthält der Text einen Syntaxfehler, zeigt der Plan den letzten gültigen
Stand und lässt sich so lange nicht ändern, bis der Fehler behoben ist.

\section{Format des Abschnitts \texttt{schematic}}

Der Plan steht in der Datei unter \code{"schematic"} – auf oberster Ebene
für die Hauptschaltung und in jeder Teilschaltung für deren eigenen Plan:

\begin{json}
"schematic": {
  "components": { "R1": { "at": [4, -2], "rot": 270, "mirror": true } },
  "wires":   [ [[0, -2], [2, -2]] ],
  "labels":  [ { "name": "out", "at": [10, -2], "rot": 0 } ],
  "ports":   [ { "name": "in", "at": [-4, 0] } ],      // nur in Teilschaltungen
  "grounds": [ { "at": [0, 2] } ]
}
\end{json}

Koordinaten sind ganze Rasterpunkte, $y$ zeigt nach unten. \code{rot} ist
die Drehung in Grad im Uhrzeigersinn (0, 90, 180, 270), \code{mirror}
spiegelt. Eine Leitung ist eine Liste von Punkten. Man muss diesen Abschnitt
nicht von Hand schreiben, kann es aber; der Simulator überliest ihn.
